% !TEX root = ../algebra_02.tex

\section{Пространство линейных отображений. Изоморфизм с пространством матриц}

\begin{Definition}{Пространство линейных отображений}{}
	Множество всех линейных отображений из $V$ в $W$ обозначается через $\Hom(V,W)$. Оно является линейным пространством относительно операций поточечного сложения и умножения на скаляр:
	\[
	(\alpha+\beta)(v)=\alpha v+\beta v, \qquad (\lambda\alpha)(v)=\lambda\alpha v
	\].
\end{Definition}

Существует естественная биекция между пространством линейных отображений и пространством матриц при фиксированных базисах.

\begin{Theorem}{Матрицы как координаты линейных отображений}{}
	Пусть $E$ --- базис $V$, $F$ --- базис $W$. Отображение
	\[
	\Phi\colon \Hom(V,W)\to M_{m\times n}(K), \qquad \Phi(\alpha)=[\alpha]_{E,F},
	\]
	является изоморфизмом линейных пространств.
\end{Theorem}

\begin{proof}
	Линейность отображения $\Phi$ напрямую следует из свойств матриц линейных отображений: $[\alpha+\beta]_{E,F}=[\alpha]_{E,F}+[\beta]_{E,F}$ и $[\lambda\alpha]_{E,F}=\lambda[\alpha]_{E,F}$.

	Покажем, что отображение биективно. Матрица $[\alpha]_{E,F}$ полностью определяется векторами образов базиса $\alpha e_1,\ldots,\alpha e_n$. Обратно, если дана произвольная матрица $A=(a_1,\ldots,a_n)$ со столбцами $a_i\in K^m$, то можно единственным образом задать линейное отображение $\alpha$ формулами $\alpha e_i=Fa_i$. В этом случае $[\alpha]_{E,F}=A$, что доказывает сюръективность и инъективность.
\end{proof}

\begin{Corollary}{Размерность пространства линейных отображений}{}
	Поскольку пространства изоморфны, их размерности совпадают. Если $\dim V=n$ и $\dim W=m$, то:
	\[
	\dim \Hom(V,W)=mn
	\].
\end{Corollary}
